\ncertchapter{3}{The Law of Radioactive Decay}\chlabel{ch:decaylaw}

\chref{ch:decay} described \emph{what} an unstable nucleus emits and how much energy the emission
carries. This chapter asks \emph{when}. The answer is surprising twice over. For a single nucleus there
is no answer at all: nothing in physics predicts the moment at which a particular nucleus will decay.
For a sample of many nuclei the answer is a single exponential law, as smooth and as reliable as any
law in mechanics, and it governs every timing question in the rest of this volume, from the activity of
a hospital source to the age of the Earth.

The chapter builds that law from one assumption, a fixed probability of decay per unit time, and then
applies it in widening circles. It defines the three clocks of a decay (the decay constant, the
half-life and the mean life) and the activity that a detector actually measures. It then treats
samples that are not a single isotope left alone: mixtures of independent isotopes, nuclei that can
decay in more than one way, chains in which the daughter is itself radioactive, and isotopes made at a
steady rate in a reactor. The mathematics is the first-order rate equation of \chref{ch:toolkit}
(Section~\ref{sec:explog}); here the work is the physics that leads to it, and the question of where it
stops being true.

A word on notation. In this volume the neutron number is always written $A-Z$, so that $N$ is free for
its universal meaning here: $N$ or $N(t)$ is the number of \emph{undecayed} nuclei in a sample at time
$t$, a pure number.

\scopebox{The decay law, half-life, mean life and activity are examined in JEE Advanced; successive
decay, equilibrium and activation are a regular source of its harder problems. The same mathematics is
the first-order kinetics of the chemistry syllabus in every examination, so the whole chapter repays
study even where the physics syllabus has been trimmed.}

\section{A Law with No Memory}\label{sec:decaylaw-memory}

The first thing to settle is what can be predicted at all. A uranium-238 nucleus formed in a supernova
five billion years ago is exactly as likely to decay in the next second as one formed yesterday.
Nuclei do not age: an old nucleus has no worn parts and no accumulated damage, because the process that
ends its life is a quantum transition whose probability depends only on the nucleus's present
structure, not on its history. That is the one physical assumption of this chapter.

\begin{defbox}{Decay constant}
Each nucleus of a given radioactive isotope has the same, constant probability per unit time of
decaying, the \textbf{decay constant} $\lambda$. In a short interval $dt$ (short enough that
$\lambda\,dt\ll1$),
\[
\Pr\big(\text{a given nucleus decays between }t\text{ and }t+dt\;\big|\;\text{it has survived to }t\big)=\lambda\,dt .
\]
$\lambda$ is a scalar with units of inverse time (\si{\per\second}, or \si{\per\year} for long-lived
isotopes). It is a rate, not a probability; only the product $\lambda\,dt$ for a short interval is a
probability. It does not depend on how long the nucleus has existed, and (to a very good
approximation, discussed below) not on temperature, pressure or chemical state.
\end{defbox}

The constancy of $\lambda$ is stated here without proof, because the reason lies in quantum
mechanics, beyond this book: it is an experimental fact, checked for many isotopes over many orders of
magnitude of time. \chref{ch:toolkit} (Section~\ref{sec:averages}) showed what such a rule implies for a
single system: the probability of surviving to time $t$ is $e^{-\lambda t}$, and the average waiting
time is $1/\lambda$. For one nucleus that is all that can be said. A sample, however, contains an
enormous number of nuclei (about \num{2.7e21} in a gram of radium), and Figure~\ref{fig:decaylaw-randomgrid}
shows what happens when many such unpredictable events are added up.

\begin{figure}[htbp]
\centering
\tikzfig{decaylaw_randomgrid}
\figcaption{fig:decaylaw-randomgrid}{Four hundred identical nuclei, each given the same chance of
decaying per unit time (filled: undecayed; open: decayed). Which nuclei decay is completely
unpredictable; how many survive is not: close to half after each half-life ($206$, $110$, $52$ here,
against an exact $200$, $100$, $50$). The scatter is about $\sqrt{100}=10$ nuclei at the first
half-life; with $10^{20}$ nuclei instead of $400$ the relative fluctuations become far too small to
measure. (The decay times were drawn once with a fixed pseudo-random sequence, so the pattern is
reproducible.)}
\end{figure}

\begin{derivbox}{The exponential decay law}
Let $N(t)$ nuclei remain undecayed at time $t$. In the next interval $dt$, each of them decays with
probability $\lambda\,dt$, independently of the others, so the expected number of decays is
$N\lambda\,dt$. The population falls by that amount:
\[
dN=-\lambda N\,dt\qquad\Longrightarrow\qquad \frac{dN}{dt}=-\lambda N .
\]
This is the first-order decay equation of \chref{ch:toolkit} (Section~\ref{sec:explog}). Separating
the variables and integrating from $N_0$ undecayed nuclei at $t=0$,
\[
\int_{N_0}^{N}\frac{dN'}{N'}=-\lambda\int_0^t dt'\quad\Longrightarrow\quad\ln\frac{N}{N_0}=-\lambda t
\quad\Longrightarrow\quad N(t)=N_0\,e^{-\lambda t}.
\]
Two checks: at $t=0$ this gives $N_0$, and as $t\to\infty$ it tends to zero without ever reaching it,
as a population that loses a fixed \emph{fraction} per unit time must.
\end{derivbox}

\begin{keybox}{The law of radioactive decay}
For a sample of a single radioactive isotope, with nothing feeding it and nothing removed except by
decay,
\begin{equation}\label{eq:decaylaw-exp}
\frac{dN}{dt}=-\lambda N,\qquad N(t)=N_0\,e^{-\lambda t}.
\end{equation}
The \emph{fractional} loss per unit time, $-(dN/dt)/N$, is the constant $\lambda$.
\end{keybox}

That last sentence is worth reading as a general statement. Any quantity that loses a fixed fraction of
itself per unit time decays exponentially, whatever the mechanism: the charge on a capacitor
discharging through a resistor, the intensity of light inside an absorbing slab, the concentration of a
drug in the blood, and the reactant in a first-order chemical reaction
(Section~\ref{sec:decaylaw-kinetics}). Radioactive decay is the cleanest case, because the fraction
lost is set by the nucleus alone.

\inwords{Each nucleus is a die rolled afresh every instant; no nucleus knows how long it has waited.
Because the same fraction of the survivors is lost in every equal interval, the number left falls by the
same \emph{factor} in every equal interval, and that is what an exponential is.}

\begin{trapbox}{the smooth law is a statement about large numbers}
Equation~\eqref{eq:decaylaw-exp} gives the \emph{expected} number of survivors. For a handful of nuclei it
says nothing useful about any one of them, and even for $400$ nuclei the actual count scatters about the
curve (Figure~\ref{fig:decaylaw-randomgrid}). A count of $n$ random events has a typical scatter of about
$\sqrt n$ (Section~\ref{sec:decaylaw-counting}), so the relative fluctuation $1/\sqrt n$ is negligible
only when $n$ is large. ``After one half-life exactly half remain'' is false for any finite sample; it is
the average over many repetitions.
\end{trapbox}

\subsection{Where the law stops working}
The derivation assumed three things, and each has a limit.
First, a large number of nuclei, as the caution above explains. Second, a closed sample: nothing
produced and nothing removed except by decay. When nuclei are fed in, by a parent's decay or by a
reactor, the equation gains a source term; Sections~\ref{sec:decaylaw-chains}
and~\ref{sec:decaylaw-activation} treat exactly that. Third, a $\lambda$ that is truly constant. Heat,
pressure and chemistry change nuclear decay rates by at most a fraction of a percent, because nuclear
energies (MeV) exceed thermal and chemical energies (a few eV at most) by a factor of about a million;
the one measurable effect is in electron capture, whose rate depends on how many electrons sit at the
nucleus. Strip an atom of \emph{all} its electrons and the change can be dramatic: a bare
rhenium-187 nucleus, which as a neutral atom has a half-life of about $42$ billion years, decays with a
half-life of about $33$ years, because a decay channel closed to the neutral atom (placing the emitted
electron directly into an empty inner orbit) opens. Such conditions exist only in storage rings and
stellar interiors. Finally, quantum mechanics predicts tiny departures from the exponential at extremely
short and extremely long times, far outside anything measured for a nucleus. For every problem in this
book, $\lambda$ is a fixed property of the isotope.

\section{Half-Life and Mean Life}\label{sec:decaylaw-halflife}

The decay constant is the natural parameter of the theory, but it is not the number most people
remember. A clock is more vivid when it is stated as a time, and two times are in use. The first answers
``how long until half is gone?''

\begin{defbox}{Half-life}
The \textbf{half-life} $\Thalf$ of an isotope is the time in which half of any sample of it decays.
Setting $N=N_0/2$ in Eq.~\eqref{eq:decaylaw-exp} gives $e^{-\lambda\Thalf}=\tfrac12$, so
\begin{equation}\label{eq:decaylaw-half}
\Thalf=\frac{\ln2}{\lambda}=\frac{0.693}{\lambda}.
\end{equation}
After $n$ half-lives a fraction $(1/2)^n$ remains, so equivalently $N=N_0\,2^{-t/\Thalf}$, a form
that suits mental arithmetic.
\end{defbox}

The half-life does not depend on the size of the sample: a gram and a tonne of radium both lose half
their nuclei in the same \SI{1600}{\year}. That independence of $N_0$ is the fingerprint of a
first-order process, and it is what Figure~\ref{fig:decaylaw-expdecay} shows: the curve halves in every
half-life, wherever on the curve one starts. On a logarithmic vertical axis,
$\ln N=\ln N_0-\lambda t$, so the decay becomes a straight line whose slope is $-\lambda$; this is how
half-lives are read from data (Section~\ref{sec:decaylaw-counting}).

\begin{figure}[htbp]
\centering
\tikzfig{decaylaw_expdecay}
\figcaption{fig:decaylaw-expdecay}{Exponential decay. Left: the population halves in every half-life,
whatever its starting size; the mean life $\tau$ is where $N$ has fallen to $N_0/e$, and the dashed
rectangle of height $N_0$ and width $\tau$ has the same area as the shaded region under the curve.
Right: on a logarithmic vertical axis the same decay is a straight line, falling by the same factor in
every half-life; its slope gives $\lambda$.}
\end{figure}

Individual nuclei live for very different times: some decay at once, some last many half-lives. Their
\emph{average} lifetime is the second clock, the mean life.

\begin{derivbox}{Mean life $\tau=1/\lambda$}
This is the mean waiting time of \chref{ch:toolkit} (Section~\ref{sec:averages}), recast for a
population of nuclei. Of $N_0$ nuclei present at $t=0$, the number that decay between $t$ and $t+dt$ is
$\lambda N\,dt=\lambda N_0e^{-\lambda t}\,dt$, and each of these lived for a time $t$. The fraction of
nuclei with lifetimes in that interval is therefore $\lambda e^{-\lambda t}\,dt$, a probability
distribution of lifetimes, and the average lifetime is the average of $t$ over it
(\chref{ch:toolkit}, Section~\ref{sec:averages}):
\[
\tau=\int_0^\infty t\,\lambda e^{-\lambda t}\,dt=\lambda\cdot\frac{1}{\lambda^2}=\frac1\lambda ,
\]
using integration by parts, $\int_0^\infty te^{-\lambda t}\,dt=\Big[-\dfrac{t\,e^{-\lambda t}}{\lambda}\Big]_0^\infty
+\dfrac1\lambda\displaystyle\int_0^\infty e^{-\lambda t}\,dt=0+\dfrac{1}{\lambda^2}$. Hence
\begin{equation}\label{eq:decaylaw-tau}
\tau=\frac1\lambda=\frac{\Thalf}{\ln2}=1.443\,\Thalf .
\end{equation}
At $t=\tau$, $N=N_0/e=0.368\,N_0$. Geometrically, the area under the decay curve is
$\int_0^\infty N_0e^{-\lambda t}\,dt=N_0/\lambda=N_0\tau$, so $\tau$ is the width of a rectangle of height $N_0$
with the same area (Figure~\ref{fig:decaylaw-expdecay}). Check of dimensions: $\lambda$ is an inverse
time, so $1/\lambda$ is a time.
\end{derivbox}

\begin{keybox}{Three clocks, one law}
\[
\lambda=\frac{\ln2}{\Thalf}=\frac1\tau,\qquad
N(t)=N_0\,e^{-\lambda t}=N_0\,e^{-t/\tau}=N_0\,2^{-t/\Thalf}.
\]
The probability that one nucleus survives to time $t$ is $e^{-\lambda t}$; the probability that it
decays between $t_1$ and $t_2$ is $e^{-\lambda t_1}-e^{-\lambda t_2}$.
\end{keybox}

\begin{trapbox}{half-life is not half the lifetime}
A sample does not ``finish'' decaying after two half-lives: after two half-lives one quarter still
remains, and after ten about $0.1\%$. Nor is the mean life equal to the half-life: it is $44\%$
longer. The few very long-lived survivors pull the average up, just as a few billionaires raise a
country's mean income above its median. The half-life is in fact the \emph{median} lifetime: half the
nuclei live less than $\Thalf$, half live longer.
\end{trapbox}

\begin{exbox}{One law written three ways}
Check that the three forms of $N(t)$ agree at $t=3\Thalf$.

\stepn{1} $N_0\,2^{-t/\Thalf}=N_0\,2^{-3}=N_0/8$.

\stepn{2} $N_0e^{-\lambda t}$ with $\lambda t=(\ln2/\Thalf)(3\Thalf)=3\ln2$: $e^{-3\ln2}=(e^{\ln2})^{-3}=2^{-3}=1/8$.

\stepn{3} $N_0e^{-t/\tau}$ with $t/\tau=3\Thalf/(\Thalf/\ln2)=3\ln2$: the same $1/8$.

The three forms are one function; choose whichever makes the arithmetic easiest.
\end{exbox}

\begin{exbox}{Fraction remaining and elapsed time: iodine-131}
Iodine-131 ($\Thalf=\SI{8.02}{\day}$) is used to treat thyroid disease. What fraction of a dose
remains after \SI{40.1}{\day}? How long until only $10\%$ remains?

\stepn{1} $40.1/8.02=5$ half-lives, so $(1/2)^5=1/32$ of the dose remains (about $3\%$).

\stepn{2} For $10\%$, $2^{-t/\Thalf}=0.10$ gives $t=\Thalf\log_2 10=3.32\,\Thalf=\SI{26.6}{\day}$.
The time is not a whole number of half-lives, so the logarithm is needed; $t=\ln(N_0/N)/\lambda$ gives
the same result.
\end{exbox}

\begin{exbox}{From half-life to decay constant: radium-226}
Radium-226 has $\Thalf=\SI{1600}{\year}$. Find $\lambda$ in \si{\per\second} and the mean life.

\stepn{1} With $\SI{1}{\year}=\SI{3.156e7}{s}$, $\lambda=\dfrac{0.693}{1600\times\SI{3.156e7}{s}}=\SI{1.37e-11}{\per\second}$.

\stepn{2} $\tau=\Thalf/\ln2=1600/0.693=\SI{2308}{\year}$, about $2300$ years.
\end{exbox}

Memorylessness has a consequence that surprises almost everyone. Because $\lambda$ does not depend on
age, a nucleus that has already survived for a time $t_1$ has exactly the same future as a freshly
made one: its probability of surviving a further time $t$ is still $e^{-\lambda t}$, and its expected
remaining life is still $\tau$. A nucleus is never ``due'' to decay.

\section{Activity}\label{sec:decaylaw-activity}

Detectors count decays, not nuclei. A Geiger counter clicks once for each emitted particle it catches,
so the quantity that experiments measure directly is the rate of decay. It deserves its own name and
symbol.

\begin{defbox}{Activity}
The \textbf{activity} $\Act$ of a sample is the number of decays per unit time:
\begin{equation}\label{eq:decaylaw-activity}
\Act=-\frac{dN}{dt}=\lambda N=\Act_0\,e^{-\lambda t},
\end{equation}
where $\Act_0=\lambda N_0$. Its SI unit is the \textbf{becquerel}: $\SI{1}{\becquerel}=1$ decay per
second. The older unit, the \textbf{curie}, is $\SI{1}{\curie}=\SI{3.7e10}{\becquerel}$, roughly the
activity of one gram of radium-226. The activity per unit mass, in \si{\becquerel\per\gram}, is the
\textbf{specific activity}; for a pure isotope of molar mass $M_\text{molar}$ it is
$\lambda N_A/M_\text{molar}$, where $N_A=\SI{6.022e23}{\per\mole}$ is Avogadro's number.
\end{defbox}

Activity decays with the same half-life as $N$, because it is $N$ multiplied by a constant. But
activity is not the number of nuclei: to go from a measured activity to $N$ one must also know
$\lambda$. The factor $\lambda$ makes the same number of nuclei very different in practice.

\begin{keybox}{Short life, high activity}
For a fixed number of nuclei, $\Act=\lambda N=N\ln2/\Thalf$: activity is inversely proportional to the
half-life. Short-lived isotopes are fiercely active but soon gone; long-lived ones are mild but
persistent.
\end{keybox}

The contrast is large. One gram of polonium-210 ($\Thalf=\SI{138.4}{\day}$, molar mass
\SI{210}{g/mol}) has slightly more nuclei than one gram of radium-226 ($\SI{1600}{\year}$,
\SI{226}{g/mol}), but a half-life $4200$ times shorter, so it emits about $4500$ times as many $\alpha$
particles per second: about \SI{1.7e14}{\becquerel}. Each decay releases \SI{5.41}{\MeV}, so the gram
gives off about \SI{140}{W} of heat, enough to raise it to several hundred degrees, and the air around it
glows faintly blue from ionisation.

\begin{exbox}{Activity and mass: the origin of the curie}
Find the activity of \SI{1.00}{g} of radium-226 ($\Thalf=\SI{1600}{\year}$, molar mass
\SI{226.0}{g/mol}).

\stepn{1} Number of nuclei: $N=\dfrac{1.00}{226.0}\times\num{6.022e23}=\num{2.66e21}$.

\stepn{2} Decay constant, in \si{\per\second} (convert years to seconds early):
$\lambda=\SI{1.373e-11}{\per\second}$ from the previous example.

\stepn{3} $\Act=\lambda N=(\num{1.373e-11})(\num{2.66e21})=\SI{3.66e10}{\becquerel}\approx\SI{0.99}{\curie}$.

The curie was defined from early measurements of exactly this sample; the small difference reflects
improved half-life data since.
\end{exbox}

\begin{remarkbox}{The smoke detector on the ceiling}
Most household ionisation smoke detectors contain about \SI{0.3}{\micro\gram} of americium-241
($\Thalf=\SI{432}{\year}$), an activity of roughly \SI{37}{\kilo\becquerel} (one microcurie). Its
$\alpha$ particles ionise the air between two plates, so a small current flows. Smoke particles soak up
the ions, the current drops, and the alarm sounds. A half-life of centuries means the source never needs
replacing, and $\alpha$ particles cannot pass through the plastic case.
\end{remarkbox}

\section{Measuring Activity in Practice}\label{sec:decaylaw-counting}

The activity is defined as decays per second, but no detector registers every decay, and some of what
it registers is not from the sample at all. Four corrections separate the true activity from the raw
count rate, and they matter as soon as a problem gives a detector reading rather than an activity.

\begin{enumerate}[label=\textbf{(\arabic*)}]
\item \textbf{Geometry.} A small source emits equally in all directions. A detector whose window
subtends a solid angle $\Omega$ at the source intercepts only the fraction $\Omega/4\pi$ of the
radiation; for a window of area $a$ at a distance $r$ large compared with its size, $\Omega\approx a/r^2$.
\item \textbf{Intrinsic efficiency} $\varepsilon$. Of the radiation entering the detector, only the
fraction $\varepsilon$ produces a count; a gamma ray, for instance, may pass straight through without
interacting.
\item \textbf{Background.} Cosmic rays and natural radioactivity in the surroundings produce counts even
with no source present. This background rate, measured separately with the source removed, must be
subtracted \emph{before} any other analysis.
\item \textbf{Statistical uncertainty.} Decay is random, so repeated counts over equal times scatter. A
count of $n$ events has an uncertainty of about $\sqrt n$ (stated here without proof; it follows from
the statistics of independent random events, and \chref{ch:nucapps} demonstrates it by simulation).
Counting $100$ events gives a precision of $10\%$; $10\,000$ events give $1\%$.
\end{enumerate}

\begin{keybox}{From count rate to activity}
\[
\text{net count rate}=\text{raw rate}-\text{background}=\Act\times\frac{\Omega}{4\pi}\times\varepsilon .
\]
The number of counts in a time $t$ is the net rate times $t$, provided $t$ is short compared with
$\Thalf$ (so the activity hardly changes); otherwise use the number of decays in the interval, below.
\end{keybox}

How many decays happen in an interval? Between $t_1$ and $t_2$ the number of nuclei falls from
$N_0e^{-\lambda t_1}$ to $N_0e^{-\lambda t_2}$, and every nucleus lost has decayed, so
\[
\text{decays between }t_1\text{ and }t_2=N_0\big(e^{-\lambda t_1}-e^{-\lambda t_2}\big),\qquad
\text{decays ever}=N_0=\frac{\Act_0}{\lambda}=\Act_0\,\tau .
\]
The second result says that a sample's total output is its initial activity sustained for one mean
life: another meaning of $\tau$.

\begin{exbox}{Decays in an interval}
A sample of \num{1.0e6} nuclei has $\Thalf=\SI{10}{min}$. How many decay between $t=\SI{10}{min}$
and $t=\SI{20}{min}$?

\stepn{1} The times are whole half-lives, so use $2^{-t/\Thalf}$: $N_0\big(2^{-1}-2^{-2}\big)=\num{1.0e6}\times(0.50-0.25)=\num{2.5e5}$.

Using the initial activity instead, $\Act_0\times\SI{10}{min}=\lambda N_0\times\SI{600}{s}=\num{6.9e5}$,
would be badly wrong: the activity falls by half within each interval.
\end{exbox}

The counts collected at successive times are also how a half-life is measured. After background is
subtracted, the net rate either shows equal-ratio intervals that can be read by eye, or is plotted as
$\ln(\text{rate})$ against $t$, which is a straight line of slope $-\lambda$.

\begin{exbox}{Extracting a half-life from measured data}
A detector records the following net count rates.

\smallskip
\begin{center}\begin{tabular}{lccccc}\toprule
$t$ (\si{min}) & 0 & 10 & 20 & 30 & 40\\
rate (\si{\per\second}) & 800 & 566 & 400 & 283 & 200\\\bottomrule\end{tabular}\end{center}

Find $\Thalf$ and $\lambda$.

\stepn{1} The rate halves every \SI{20}{min} ($800\to400\to200$), so $\Thalf=\SI{20}{min}$ and
$\lambda=0.693/20=\SI{0.0347}{\per\minute}$.

\stepn{2} Check the intermediate points: in \SI{10}{min} the factor is $2^{-1/2}=0.707$, and
$800\times0.707=566$, $400\times0.707=283$. On a plot of $\ln(\text{rate})$ against $t$ the points lie
on a line of slope \SI{-0.0347}{\per\minute}.
\end{exbox}

\begin{trapbox}{subtract the background first}
A constant background added to a decaying signal makes the late counts fall too slowly, so the
successive ratios drift and the half-life appears to lengthen with time. A decay that looks
``non-exponential'' in raw data is often a pure exponential plus an unsubtracted background.
\end{trapbox}

\section{Mixtures of Independent Isotopes}\label{sec:decaylaw-mixtures}

A real sample is often not one isotope but several that have nothing to do with each other: two fission
products, say, or two isotopes made together in a reactor. They are labelled $1$ and $2$ here (these
subscripts always mean \emph{independent} isotopes; the letters X and Y are kept for parent and
daughter). Each decays by its own law, unaffected by the other, so the total activity is simply the sum:
\[
\Act_\text{total}(t)=\Act_{1,0}\,e^{-\lambda_1t}+\Act_{2,0}\,e^{-\lambda_2t},\qquad
\frac{N_1}{N_2}=\frac{N_{1,0}}{N_{2,0}}\,e^{-(\lambda_1-\lambda_2)t}.
\]
On a semi-log plot this sum is \emph{not} a straight line. It is curved, steep at first and flatter
later (Figure~\ref{fig:decaylaw-peeling}), because the short-lived isotope dominates early but soon dies
away, leaving only the long-lived one. That suggests a way to take a measured curve apart.

\begin{methodbox}{Curve peeling}
\begin{enumerate}[label=\arabic*.,leftmargin=*,itemsep=1pt]
\item Fit a straight line to the late-time tail of $\ln\Act$ against $t$; its slope gives the long
half-life.
\item Extrapolate that line back to $t=0$ and subtract it from the data at every time.
\item The remainder is a new straight line on semi-log axes; its slope gives the short half-life.
\end{enumerate}
\end{methodbox}

\begin{figure}[htbp]
\centering
\tikzfig{decaylaw_peeling}
\figcaption{fig:decaylaw-peeling}{Total activity of two independent isotopes (half-lives \SI{1}{h} and
\SI{10}{h}, each starting at \SI{1000}{\becquerel}) on a logarithmic axis. Notice that the total is
curved. Beyond about \SI{6}{h} the short-lived part is under $3\%$ of the total, so the tail is
straight and belongs to the long-lived isotope. Subtracting its back-extrapolation leaves the
short-lived component, another straight line.}
\end{figure}

\begin{exbox}{Ratio of two isotopes}
Two isotopes start with equal numbers of nuclei. Their half-lives are \SI{10}{min} and \SI{20}{min}.
Find the ratio of their numbers after \SI{1}{h}.

\stepn{1} \SI{60}{min} is $6$ half-lives of the first and $3$ of the second, so
$N_1/N_2=2^{-6}/2^{-3}=1/8$.

\stepn{2} The ratio of their \emph{activities} is different, because $\lambda_1=2\lambda_2$:
$\Act_1/\Act_2=2\times\tfrac18=\tfrac14$. Read the question for which ratio it wants.
\end{exbox}

\section{Parallel Decay: Branching}\label{sec:decaylaw-branching}

So far each nucleus has had one way to decay. Some have two or more. Bismuth-212, for instance, decays
both by $\alpha$ emission and by $\beta^-$ emission (the modes of \chref{ch:decay}). Each route, or
\emph{channel}, labelled $i$, has its own probability per unit time $\lambda_i$, called a \textbf{partial
decay constant}. A nucleus that can leave by either of two doors leaves sooner than one with a single
door, and the question is how much sooner.

\begin{derivbox}{Total decay constant and branching ratio}
In a short time $dt$ the probability that a nucleus decays by channel $i$ is $\lambda_i\,dt$. The
channels are mutually exclusive (a nucleus decays only once), so the probability that it decays by
\emph{some} channel is the sum, $\big(\sum_i\lambda_i\big)dt$. The population therefore obeys
\[
\frac{dN}{dt}=-\Big(\sum_i\lambda_i\Big)N\qquad\Longrightarrow\qquad\lambda=\sum_i\lambda_i .
\]
Of the $\lambda N\,dt$ decays in $dt$, $\lambda_iN\,dt$ go through channel $i$, so the fraction going
through channel $i$, the \textbf{branching ratio}, is $f_i=\lambda_i/\lambda$, the same at every time.
Writing each $\lambda_i=\ln2/T_i$ defines the \textbf{partial half-life} $T_i$ of channel $i$, and
dividing $\lambda=\sum\lambda_i$ by $\ln2$ gives the half-life relation below.
\end{derivbox}

\begin{keybox}{Branching}
\[
\lambda=\sum_i\lambda_i,\qquad \frac{1}{\Thalf}=\sum_i\frac{1}{T_i},\qquad
f_i=\frac{\lambda_i}{\lambda},\qquad T_i=\frac{\Thalf}{f_i}.
\]
For two channels, $1/\Thalf=1/T_\alpha+1/T_\beta$, which combines like resistors in parallel: the
combination is always shorter than either part. The number of decays by channel $i$ in any interval is
$f_i$ times the total number of decays in that interval.
\end{keybox}

Two limits check the result. If one channel is switched off ($\lambda_2\to0$, so $T_2\to\infty$),
$\Thalf\to T_1$, as it must. If the two channels are equally likely, $T_1=T_2$ and $\Thalf=T_1/2$:
two equal doors halve the waiting time.

\begin{trapbox}{partial half-lives are not observed half-lives}
A partial half-life is \emph{not} the time for half the sample to decay by that route. It is the
half-life the sample \emph{would} have if every other channel were switched off. In reality the parent
population falls with the total $\lambda$, so the products of every channel grow with the same overall
half-life $\Thalf$, shorter than any partial half-life. Measuring the half-life by watching only the
$\alpha$ particles still gives $\Thalf$, not $T_\alpha$.
\end{trapbox}

\begin{exbox}{Two decay channels}
A nucleus decays by $\alpha$ emission with partial half-life \SI{1620}{\year} and by $\beta^-$ emission
with partial half-life \SI{405}{\year}. Find the half-life, the time for $3/4$ of a sample to decay,
and the fraction of decays that are $\alpha$.

\stepn{1} $1/\Thalf=1/1620+1/405=1/1620+4/1620=5/1620$, so $\Thalf=\SI{324}{\year}$, shorter than
either partial half-life.

\stepn{2} Three-quarters decays in two half-lives, \SI{648}{\year}.

\stepn{3} $f_\alpha=\lambda_\alpha/\lambda=\Thalf/T_\alpha=324/1620=0.20$, so $20\%$ of decays are
$\alpha$ and $80\%$ are $\beta^-$.
\end{exbox}

\section{Successive Decay: Radioactive Chains}\label{sec:decaylaw-chains}

\chref{ch:decay} showed that a heavy nucleus rarely reaches stability in one step: uranium-238 passes
through fourteen decays before ending as lead-206. Every intermediate member is both produced (by its
parent's decay) and destroyed (by its own). The simplest chain has two radioactive members: a parent X
decaying with constant $\lambda_X$ into a daughter Y, which decays with constant $\lambda_Y$ into a
stable nuclide,
$\mathrm{X}\xrightarrow{\;\lambda_X\;}\mathrm{Y}\xrightarrow{\;\lambda_Y\;}\text{stable}$, with
half-lives $T_X$ and $T_Y$ and populations $N_X$ and $N_Y$. Each decay of X makes one Y, so the
daughter is fed at the rate $\lambda_XN_X$ and drained at the rate $\lambda_YN_Y$:
\begin{equation}\label{eq:decaylaw-chain}
\frac{dN_X}{dt}=-\lambda_XN_X,\qquad\frac{dN_Y}{dt}=\lambda_XN_X-\lambda_YN_Y .
\end{equation}
The first equation is the ordinary decay law; the second is the decay law with a source term, and its
source is itself decaying.

\begin{derivbox}{The daughter population (Bateman solution)}
Start with pure parent: $N_X(0)=N_0$ and $N_Y(0)=0$. Then $N_X=N_0e^{-\lambda_Xt}$, and the daughter's
equation becomes
\[
\frac{dN_Y}{dt}+\lambda_YN_Y=\lambda_XN_0e^{-\lambda_Xt}.
\]
Multiply both sides by $e^{\lambda_Yt}$. By the product rule the left side is exactly
$\frac{d}{dt}\big(N_Ye^{\lambda_Yt}\big)$ (this multiplier is called an \emph{integrating factor}), so
\[
\frac{d}{dt}\big(N_Ye^{\lambda_Yt}\big)=\lambda_XN_0e^{(\lambda_Y-\lambda_X)t}
\;\Longrightarrow\;
N_Ye^{\lambda_Yt}=\frac{\lambda_XN_0}{\lambda_Y-\lambda_X}\big(e^{(\lambda_Y-\lambda_X)t}-1\big),
\]
where the constant of integration was fixed by $N_Y(0)=0$. Therefore
\begin{equation}\label{eq:decaylaw-bateman}
N_Y(t)=\frac{\lambda_X N_0}{\lambda_Y-\lambda_X}\big(e^{-\lambda_Xt}-e^{-\lambda_Yt}\big),\qquad
\Act_Y(t)=\lambda_YN_Y=\Act_X(0)\,\frac{\lambda_Y}{\lambda_Y-\lambda_X}\big(e^{-\lambda_Xt}-e^{-\lambda_Yt}\big).
\end{equation}
$N_Y$ starts at zero, rises and eventually falls. At its maximum $dN_Y/dt=0$; differentiating
Eq.~\eqref{eq:decaylaw-bateman} gives $-\lambda_Xe^{-\lambda_Xt}+\lambda_Ye^{-\lambda_Yt}=0$, so
$e^{(\lambda_Y-\lambda_X)t}=\lambda_Y/\lambda_X$ and the maximum occurs at
\begin{equation}\label{eq:decaylaw-tmax}
t_{\max}=\frac{\ln(\lambda_Y/\lambda_X)}{\lambda_Y-\lambda_X}.
\end{equation}
At that instant the second of Eqs.~\eqref{eq:decaylaw-chain} reads $0=\lambda_XN_X-\lambda_YN_Y$: the
parent and daughter activities are momentarily equal.
\end{derivbox}

The formula looks as if it fails when $\lambda_Y=\lambda_X$, where both numerator and denominator
vanish, so it is worth checking that limit. Write $\lambda_Y=\lambda_X+\delta$ with $\delta$ small:
$e^{-\lambda_Xt}-e^{-\lambda_Yt}=e^{-\lambda_Xt}\big(1-e^{-\delta t}\big)\approx e^{-\lambda_Xt}\,\delta t$, so
$N_Y\to\lambda N_0\,t\,e^{-\lambda t}$, a perfectly finite curve with its maximum at $t_{\max}=1/\lambda$
(which is also the limit of Eq.~\eqref{eq:decaylaw-tmax}). The opposite limit, a very short-lived
daughter ($\lambda_Y\gg\lambda_X$), gives $N_Y\approx(\lambda_X/\lambda_Y)N_X$ once $e^{-\lambda_Yt}$ has
died away: the daughter decays as fast as it is made. Which of these behaviours appears depends on which
nucleus lives longer, and three regimes follow (Figure~\ref{fig:decaylaw-successive}).

\begin{figure}[htbp]
\centering
\tikzfig{decaylaw_successive}
\figcaption{fig:decaylaw-successive}{Activities in a two-member chain starting from pure parent
(blue: parent X; orange: daughter Y), for parent half-lives of $100$, $4$ and $\tfrac14$ times the
daughter's; time is in units of the daughter's half-life. In each panel the daughter's curve crosses the
parent's exactly at the daughter's maximum (dot). Notice that in transient equilibrium the daughter's
activity settles \emph{above} the parent's, at $\tfrac43$ of it.}
\end{figure}

\subsection{Secular equilibrium (\texorpdfstring{$T_X\gg T_Y$}{TX much longer than TY})}
After a few daughter half-lives the daughter decays as
fast as it is produced. Over any practical time $N_X$ is constant, so $dN_Y/dt=0$ and
\begin{equation}\label{eq:decaylaw-secular}
\lambda_XN_X=\lambda_YN_Y\qquad\Longleftrightarrow\qquad\frac{N_Y}{N_X}=\frac{T_Y}{T_X}.
\end{equation}
The populations are in the ratio of the half-lives: the short-lived daughter is rare. Applied link by
link, the same argument gives a striking result for a long chain in an undisturbed rock: every member
has the same activity as the long-lived head of the chain.

\subsection{Transient equilibrium (\texorpdfstring{$T_X>T_Y$}{TX longer than TY}, but not hugely)}
Since $\lambda_Y>\lambda_X$, the term
$e^{-\lambda_Yt}$ in Eq.~\eqref{eq:decaylaw-bateman} dies away first, leaving
$N_Y\approx\dfrac{\lambda_X}{\lambda_Y-\lambda_X}N_0e^{-\lambda_Xt}=\dfrac{\lambda_X}{\lambda_Y-\lambda_X}N_X$.
Multiplying by $\lambda_Y$ and dividing by $\Act_X=\lambda_XN_X$,
\begin{equation}\label{eq:decaylaw-transient}
\frac{\Act_Y}{\Act_X}=\frac{\lambda_Y}{\lambda_Y-\lambda_X}=\frac{T_X}{T_X-T_Y}>1 .
\end{equation}
The daughter then decays in step with the parent, with the \emph{parent's} half-life, and its activity
exceeds the parent's. Secular equilibrium is the limit $\lambda_X\ll\lambda_Y$, where this ratio tends
to $1$. In the other direction, as $T_X$ approaches $T_Y$ from above the ratio grows without bound, but
the time to reach the equilibrium grows too, and at $T_X=T_Y$ there is no equilibrium at all.

\subsection{No equilibrium (\texorpdfstring{$T_X<T_Y$}{TX shorter than TY})}
Now Eq.~\eqref{eq:decaylaw-transient} would give a negative
ratio, a sign that there is no steady ratio. The parent dies out first, leaving the daughter to decay
with its own half-life.

\begin{keybox}{Equilibrium in a parent--daughter chain}
\[
\text{secular } (T_X\gg T_Y):\ \Act_Y=\Act_X,\ \ \frac{N_Y}{N_X}=\frac{T_Y}{T_X};\qquad
\text{transient } (T_X>T_Y):\ \frac{\Act_Y}{\Act_X}=\frac{T_X}{T_X-T_Y}.
\]
The daughter peaks at $t_{\max}=\ln(\lambda_Y/\lambda_X)/(\lambda_Y-\lambda_X)$, where the two
activities are equal. For mass ratios, multiply number ratios by the ratio of molar masses.
\end{keybox}

\begin{remarkbox}{The bathtub picture}
Think of the daughter population as water in a bathtub. The parent's decays are the tap, pouring in at
the rate $\lambda_XN_X$. The drain empties at a rate proportional to the water level, $\lambda_YN_Y$,
because more nuclei mean more decays. Start with an empty tub: the level rises until outflow matches
inflow, and then holds steady. That steady state, $\lambda_YN_Y=\lambda_XN_X$, is secular equilibrium.
If the tap itself slowly weakens (transient equilibrium), the level follows it down, lagging slightly,
which is why the outflow then slightly exceeds the inflow. A reactor making an isotope at a constant
rate (Section~\ref{sec:decaylaw-activation}) is a tap that never weakens.
\end{remarkbox}

\begin{exbox}{Time of the daughter maximum}
A parent of half-life \SI{2.0}{h} decays to a daughter of half-life \SI{1.0}{h}. Starting with pure
parent, when is the number of daughter nuclei largest?

\stepn{1} $\lambda_X=0.693/2.0=\SI{0.347}{\per\hour}$ and $\lambda_Y=\SI{0.693}{\per\hour}$, so
$\lambda_Y/\lambda_X=2$.

\stepn{2} $t_{\max}=\dfrac{\ln2}{0.693-0.347}=\dfrac{0.693}{0.347}=\SI{2.0}{h}$. At that instant
parent and daughter activities are equal; afterwards the pair settles into transient equilibrium with
$\Act_Y/\Act_X=2/(2-1)=2$.
\end{exbox}

\begin{remarkbox}{The technetium ``cow''}
Technetium-99m ($\Thalf=\SI{6.0}{h}$) is the workhorse of medical imaging, but it cannot be stored for
long. Hospitals instead receive a shielded generator containing molybdenum-99 ($\SI{66}{h}$), which
decays to technetium-99m. Equation~\eqref{eq:decaylaw-tmax} puts the daughter's maximum at
$t_{\max}=\SI{23}{h}$, so the generator approaches transient equilibrium in about a day. Each morning a
saline wash ``milks'' the accumulated technetium off the column, and within hours it regrows. The ideal
transient ratio would be $66/(66-6.0)=1.10$; in practice about $12\%$ of molybdenum decays bypass the
metastable state, so the technetium activity settles near $0.96$ of the molybdenum's. The generator is
replaced weekly, after about $2.5$ parent half-lives.
\end{remarkbox}

\section{Production at a Constant Rate}\label{sec:decaylaw-activation}

The bathtub remark ended with a tap that never weakens, and that is what a reactor or a particle
accelerator provides. There a radioactive isotope is often \emph{made} at a constant rate $\Prod$
(nuclei per second) while it decays. Cobalt-60, the classic radiotherapy source, is made by exposing
ordinary cobalt-59 to reactor neutrons, $\iso{59}{Co}+n\to\iso{60}{Co}+\gamma$ (reactions of this kind
are the subject of \chref{ch:reactions}). This is called \textbf{activation}. The number of radioactive
nuclei then obeys
\[
\frac{dN}{dt}=\Prod-\lambda N .
\]

\begin{derivbox}{Growth to saturation}
This is the equation of growth with supply in \chref{ch:toolkit} (Section~\ref{sec:explog}). With
$N(0)=0$, the substitution $w=\Prod-\lambda N$ gives $dw/dt=-\lambda\,dN/dt=-\lambda w$, so $w$ decays
exponentially from its initial value $w(0)=\Prod$: $w=\Prod e^{-\lambda t}$. Solving for $N$,
\begin{equation}\label{eq:decaylaw-activation}
N(t)=\frac{\Prod}{\lambda}\big(1-e^{-\lambda t}\big),\qquad\Act(t)=\lambda N=\Prod\big(1-e^{-\lambda t}\big)
=\Prod\big(1-2^{-t/\Thalf}\big).
\end{equation}
Checks: at $t=0$, $N=0$; for $\lambda t\ll1$, $N\approx\Prod t$ (almost nothing has had time to decay);
as $t\to\infty$, $N\to\Prod/\lambda$ and $\Act\to\Prod$. Production and decay then balance, the same
balance as secular equilibrium, with the reactor playing the long-lived parent.
\end{derivbox}

\begin{keybox}{Activation}
Irradiation for a time $t$ gives $\Act=\Prod\big(1-2^{-t/\Thalf}\big)$: $\tfrac12$ of the saturation
activity $\Prod$ after one half-life, $\tfrac34$ after two, $\tfrac78$ after three. After the sample is
removed at $t_0$, it decays freely: $\Act=\Act(t_0)\,2^{-(t-t_0)/\Thalf}$.
\end{keybox}

Figure~\ref{fig:decaylaw-activation} shows the consequence for anyone paying for reactor time: the
approach to saturation is the decay curve turned upside down, so each further half-life of irradiation
only halves the remaining gap. Irradiating beyond three or four half-lives gains almost nothing.

\begin{figure}[htbp]
\centering
\tikzfig{decaylaw_activation}
\figcaption{fig:decaylaw-activation}{Activation by production at a constant rate $\Prod$. After one
half-life of irradiation the activity reaches $\tfrac12$ of saturation, after two $\tfrac34$, after
three $\tfrac78$: each half-life halves the remaining gap, so irradiating beyond a few half-lives gains
little. Once the source is removed (here after five half-lives), ordinary exponential decay takes
over.}
\end{figure}

\begin{exbox}{Activity during production}
A radioisotope with $\Thalf=\SI{10}{min}$ is produced at $\Prod=\num{1.0e10}$ nuclei per second. Find
its activity and the number of nuclei present after \SI{20}{min} of production.

\stepn{1} \SI{20}{min} is two half-lives, so $\Act=\Prod\big(1-\tfrac14\big)=\SI{7.5e9}{\becquerel}$.

\stepn{2} $N=\Act/\lambda$ with $\lambda=0.693/\SI{600}{s}$:
$N=\num{7.5e9}\times600/0.693=\num{6.5e12}$.
\end{exbox}

\section{A Bridge to Chemistry: First-Order Kinetics}\label{sec:decaylaw-kinetics}

The decay law is not only nuclear. A first-order chemical reaction, such as the decomposition
$\mathrm{N_2O_5}\to2\,\mathrm{NO_2}+\tfrac12\mathrm{O_2}$, obeys $d[\mathrm{R}]/dt=-k[\mathrm{R}]$, where
$[\mathrm R]$ is the reactant concentration and $k$ is the chemical rate constant. (In this section
only, $k$ is chemistry's rate constant; elsewhere in this volume $k$ is the neutron multiplication factor
of \chref{ch:reactions}.) It is the same equation with different symbols, and everything in
Sections~\ref{sec:decaylaw-memory} and~\ref{sec:decaylaw-halflife} carries over;
Table~\ref{tab:decaylaw-kinetics} sets the two side by side.

\begin{tablebox}{Radioactive decay and first-order reactions compared}
\centering
\begin{tabular}{@{}lll@{}}\toprule
 & \textbf{Radioactive decay} & \textbf{First-order reaction}\\\midrule
Quantity & number of nuclei $N$ & concentration $[\mathrm R]$\\
Rate law & $dN/dt=-\lambda N$ & $d[\mathrm R]/dt=-k[\mathrm R]$\\
Integrated form & $N=N_0e^{-\lambda t}$ & $[\mathrm R]=[\mathrm R]_0e^{-kt}$, that is,
$k=\dfrac{2.303}{t}\log_{10}\dfrac{[\mathrm R]_0}{[\mathrm R]}$\\[6pt]
Half-life & $\Thalf=0.693/\lambda$, independent of $N_0$ & $t_{1/2}=0.693/k$, independent of $[\mathrm R]_0$\\
Effect of temperature & none (nuclear energies $\gg k_BT$) & strong: $k\propto e^{-E_a/k_BT}$ (Arrhenius)\\
Microscopic picture & one nucleus, intrinsic probability & molecule must gain activation energy $E_a$\\\bottomrule
\end{tabular}\label{tab:decaylaw-kinetics}
\end{tablebox}

The last two rows explain the difference. A chemical reaction needs a thermal kick over an energy
barrier $E_a$ of order \SI{1}{eV} per molecule, not far above molecular thermal energies
($k_BT\approx\SI{0.03}{eV}$ at room temperature, with $k_B$ Boltzmann's constant), so the factor
$e^{-E_a/k_BT}$ changes enormously when the temperature is raised. The energies that govern nuclear
decay are millions of electronvolts, more than a million times $k_BT$ even in the hottest laboratory
furnace, so heat and chemistry cannot reach the nucleus. That is the physical reason $\lambda$ is a
constant of the isotope, the assumption on which this chapter was built.

\section{Summary}\label{sec:summary-decaylaw}
One assumption, a fixed probability of decay per unit time for every nucleus, yields a smooth
exponential law for large samples and, by adding sources and channels to the same equation, the
behaviour of mixtures, branching decays, chains and activated samples.

\begin{itemize}
\item $dN/dt=-\lambda N$, so $N=N_0e^{-\lambda t}=N_0e^{-t/\tau}=N_0\,2^{-t/\Thalf}$, with
      $\Thalf=\ln2/\lambda=0.693/\lambda$ and $\tau=1/\lambda=1.443\,\Thalf$. The half-life is the median
      lifetime and is independent of the sample size.
\item One nucleus survives to $t$ with probability $e^{-\lambda t}$, decays between $t_1$ and $t_2$ with
      probability $e^{-\lambda t_1}-e^{-\lambda t_2}$, and has no memory of its age.
\item Activity $\Act=\lambda N=\Act_0e^{-\lambda t}$, in becquerels, where
      $\SI{1}{\curie}=\SI{3.7e10}{\becquerel}$. For a mass $m$, $N=(m/M_\text{molar})N_A$. For a fixed $N$, activity is inversely proportional to
      half-life.
\item Net count rate $=\Act\,(\Omega/4\pi)\,\varepsilon$ after the background is subtracted; a count of
      $n$ carries an uncertainty of about $\sqrt n$. Decays between $t_1$ and $t_2$:
      $N_0(e^{-\lambda t_1}-e^{-\lambda t_2})$; decays ever: $N_0=\Act_0\tau$.
\item A semi-log plot of $\ln\Act$ against $t$ is a straight line of slope $-\lambda$ for one isotope;
      a curved one signals a mixture, taken apart by curve peeling.
\item Branching: $\lambda=\sum\lambda_i$, $1/\Thalf=\sum 1/T_i$, $f_i=\lambda_i/\lambda$, and every
      channel's products grow with the total half-life.
\item Chains: $dN_Y/dt=\lambda_XN_X-\lambda_YN_Y$, solved by the Bateman formula; the daughter peaks at
      $t_{\max}=\ln(\lambda_Y/\lambda_X)/(\lambda_Y-\lambda_X)$, where the activities are equal. Secular
      equilibrium: $\Act_Y=\Act_X$, $N_Y/N_X=T_Y/T_X$. Transient: $\Act_Y/\Act_X=T_X/(T_X-T_Y)$.
\item Production at a constant rate: $\Act=\Prod(1-e^{-\lambda t})$, saturating at $\Act=\Prod$,
      $N=\Prod/\lambda$.
\item The law fails for small samples (fluctuations), for samples that are fed or depleted (add the source
      term), and, very slightly, for electron capture in unusual chemical states; heat and chemistry do not
      change $\lambda$, which is why first-order chemical kinetics, unlike decay, depends on temperature.
\end{itemize}

\section{Problem Types}\label{sec:types-decaylaw}
Ten types cover the decay-law questions examiners set. Most of them turn on one habit: decide first
whether the question is about nuclei remaining, nuclei decayed, or decays per second, because the three
are easily confused.

\archq{Fraction remaining and elapsed time}
{``What fraction remains, or has decayed, after\dots'', ``how long until\dots''.}
{If $t$ is a whole number of half-lives, use $(1/2)^n$. Otherwise use $N/N_0=2^{-t/\Thalf}$, so
$t=\Thalf\log_2(N_0/N)$, or $t=\ln(N_0/N)/\lambda$.}
{The wording: ``decayed'' is $1-N/N_0$, not $N/N_0$.}
{Iodine-131 has a half-life of \SI{8.0}{\day}. The fraction of a sample that has \emph{decayed} after
\SI{24}{\day} is:}
{$\tfrac18$}
{$\tfrac78$}
{$\tfrac38$}
{$\tfrac34$\archans{B}{$24/8.0=3$ half-lives, so $(1/2)^3=\tfrac18$ remains and $1-\tfrac18=\tfrac78$
has decayed. Option (A) is the fraction remaining; (C) treats three half-lives as ``three eighths'';
(D) counts only two half-lives.}}

\archq{Decay constant, half-life, mean life and survival probability}
{Conversions between $\lambda$, $\Thalf$ and $\tau$; the probability that \emph{one} nucleus survives
or decays in a given window.}
{$\lambda=0.693/\Thalf=1/\tau$, so $\Thalf=0.693\,\tau$. Survival to $t$: $e^{-\lambda t}$; decay between
$t_1$ and $t_2$: $e^{-\lambda t_1}-e^{-\lambda t_2}$. A nucleus that has already survived has the same
future as a fresh one.}
{Treating the mean life as the half-life, or multiplying by $\ln2$ when one should divide; the mean life
is the \emph{longer} of the two.}
{The mean life of a radioactive isotope is \SI{10}{min}. Its half-life is:}
{\SI{14.4}{min}}
{\SI{5.0}{min}}
{\SI{6.9}{min}}
{\SI{10}{min}\archans{C}{$\Thalf=\tau\ln2=0.693\times10=\SI{6.9}{min}$. Option (A) divides by $\ln2$
instead of multiplying, which would make the half-life longer than the mean life; (B) takes the
half-life to be half the lifetime; (D) equates the mean and the median.}}

\archq{Activity and mass}
{Activity of a given mass, mass needed for a given activity, specific activity (\si{\becquerel} per
gram).}
{$N=(m/M_\text{molar})N_A$; $\lambda=0.693/\Thalf$ \emph{in} \si{\per\second}; $\Act=\lambda N$.
Convert years to seconds early. $\SI{1}{\curie}=\SI{3.7e10}{\becquerel}$.}
{Leaving $\lambda$ in \si{\per\year}, which gives ``becquerels'' too small by a factor of
\num{3.2e7}, or forgetting the $\ln2$.}
{The activity of \SI{1.0}{mg} of cobalt-60 ($\Thalf=\SI{5.27}{\year}$, molar mass \SI{60}{g/mol};
$\SI{1}{\year}=\SI{3.156e7}{s}$) is about:}
{\SI{4.2e10}{\becquerel}}
{\SI{6.0e10}{\becquerel}}
{\SI{4.2e13}{\becquerel}}
{\SI{1.3e18}{\becquerel}\archans{A}{$N=(10^{-3}/60)(\num{6.022e23})=\num{1.0e19}$ and
$\lambda=0.693/(5.27\times\num{3.156e7})=\SI{4.2e-9}{\per\second}$, so $\Act=\SI{4.2e10}{\becquerel}$
(about \SI{1.1}{\curie}). Option (B) uses $N/\Thalf$, forgetting $\ln2$; (C) uses a gram instead of a
milligram; (D) leaves $\lambda$ in \si{\per\year}, so the number is decays per year, not per second.}}

\archq{Decays in an interval; counting}
{Number of decays between two times, total decays ever, counts registered by a detector.}
{Decays between $t_1$ and $t_2$: $N_0(e^{-\lambda t_1}-e^{-\lambda t_2})$. Total ever: $N_0=\Act_0\tau$.
Counts $=$ activity $\times\,(\Omega/4\pi)\times\varepsilon\times$ time, valid only for times short
compared with $\Thalf$; add any background.}
{Multiplying the initial activity by a long time, as if the activity stayed constant.}
{A sample contains \num{8.0e6} nuclei of an isotope with $\Thalf=\SI{1.0}{h}$. The number that decay
between $t=\SI{1.0}{h}$ and $t=\SI{3.0}{h}$ is:}
{\num{1.0e6}}
{\num{4.0e6}}
{\num{7.0e6}}
{\num{3.0e6}\archans{D}{$N_0\big(2^{-1}-2^{-3}\big)=\num{8.0e6}\times\big(\tfrac12-\tfrac18\big)=\num{3.0e6}$.
Option (A) is the number still undecayed at \SI{3}{h}; (B) is the number that decay in the first hour
(or remain at \SI{1}{h}); (C) is the number that decay in the whole first three hours.}}

\archq{Extracting half-lives from measured data}
{A table or graph of count rates against time, possibly including background.}
{Subtract background first. Then spot equal-ratio intervals (the halving time), or fit a straight line
to $\ln(\text{rate})$ against $t$, whose slope is $-\lambda$.}
{Taking ratios of raw rates; or interpolating linearly between readings, which treats an exponential as
a straight line.}
{A counter reads \SI{850}{\per\second} at $t=0$ and \SI{250}{\per\second} at $t=\SI{2.0}{h}$. With
the source removed it reads \SI{50}{\per\second}. The half-life of the source is:}
{\SI{0.50}{h}}
{\SI{1.0}{h}}
{\SI{1.13}{h}}
{\SI{1.33}{h}\archans{B}{Net rates are $800$ and \SI{200}{\per\second}: a factor of $4=2^2$ in
\SI{2.0}{h}, so $\Thalf=\SI{1.0}{h}$. Option (C) uses the raw ratio $850/250=3.4$, ignoring the
background; (D) interpolates linearly ($\tfrac34$ lost in \SI{2}{h}, so half in \SI{1.33}{h}); (A)
divides the time by the factor $4$.}}

\archq{Mixtures of two isotopes}
{Two independent isotopes in one sample; ratios at a later time; when activities or numbers are equal;
a curved semi-log activity plot.}
{Each decays independently; the total activity is the sum. $N_1/N_2=(N_{1,0}/N_{2,0})e^{-(\lambda_1-\lambda_2)t}$.
For equal activities set $\lambda_1N_{1,0}e^{-\lambda_1t}=\lambda_2N_{2,0}e^{-\lambda_2t}$ and solve
for $t$. Late times belong to the longer-lived isotope.}
{Confusing a ratio of numbers with a ratio of activities: they differ by $\lambda_1/\lambda_2$.}
{Isotopes 1 ($\Thalf=\SI{1.0}{h}$) and 2 ($\Thalf=\SI{3.0}{h}$) have equal \emph{activities} at
$t=0$. After \SI{6.0}{h}, the ratio of their \emph{numbers of nuclei} $N_1/N_2$ is:}
{$\tfrac1{16}$}
{$\tfrac3{16}$}
{$\tfrac1{48}$}
{$\tfrac13$\archans{C}{Equal activities mean $N_{1,0}/N_{2,0}=\lambda_2/\lambda_1=\tfrac13$. After
\SI{6}{h}, isotope 1 has passed $6$ half-lives and isotope 2 two, so
$N_1/N_2=\tfrac13\times2^{-6}/2^{-2}=\tfrac13\times\tfrac1{16}=\tfrac1{48}$. Option (A) is the activity
ratio; (B) multiplies by $\lambda_1/\lambda_2$ instead of dividing; (D) is the initial number ratio.}}

\archq{Branching (parallel) decay}
{One nucleus, two decay modes; partial half-lives or branching percentages.}
{$\lambda=\sum_i\lambda_i$, that is, $1/\Thalf=\sum_i1/T_i$. Branching fraction $f_i=\lambda_i/\lambda$,
so $T_i=\Thalf/f_i$. Decays by channel $i$ in any interval are $f_i$ times the total.}
{Adding partial half-lives, or taking the observed half-life to be one of them; the combination is
shorter than either.}
{A nucleus decays by $\alpha$ emission with a partial half-life of \SI{6.0}{h} and by $\beta^-$
emission with a partial half-life of \SI{3.0}{h}. The half-life of a sample is:}
{\SI{9.0}{h}}
{\SI{4.5}{h}}
{\SI{3.0}{h}}
{\SI{2.0}{h}\archans{D}{$1/\Thalf=1/6.0+1/3.0=3/6.0$, so $\Thalf=\SI{2.0}{h}$, shorter than either
partial half-life (and one-third of decays are $\alpha$). Option (A) adds the half-lives; (B) averages
them; (C) keeps only the faster channel.}}

\archq{Successive decay and equilibrium}
{Parent--daughter chains; ``time of maximum daughter''; abundance or mass ratios in old ores;
transient-equilibrium ratios.}
{Time of maximum: Eq.~\eqref{eq:decaylaw-tmax}. Secular equilibrium ($T_X\gg T_Y$):
$N_Y/N_X=T_Y/T_X$, then multiply by the ratio of molar masses for a mass ratio. Transient:
$\Act_Y/\Act_X=T_X/(T_X-T_Y)$.}
{Inverting the ratio (the short-lived daughter is the \emph{rare} one), or forgetting the molar masses
when a mass is asked for.}
{An undisturbed uranium ore has reached secular equilibrium. The mass of radium-226
($\Thalf=\SI{1600}{\year}$, \SI{226}{g/mol}) present per kilogram of uranium-238
($\Thalf=\SI{4.47e9}{\year}$, \SI{238}{g/mol}) is about:}
{\SI{0.34}{mg}}
{\SI{0.36}{mg}}
{\SI{0.38}{mg}}
{\SI{2.7e6}{kg}\archans{A}{$N_\text{Ra}/N_\text{U}=1600/\num{4.47e9}=\num{3.58e-7}$; the mass ratio is
$\num{3.58e-7}\times226/238=\num{3.40e-7}$, so \SI{0.34}{mg} per kilogram. Option (B) omits the molar
masses; (C) inverts their ratio; (D) inverts the half-life ratio, making the short-lived daughter the
abundant one.}}

\archq{Production at a constant rate (activation)}
{Nuclei produced at a steady rate $\Prod$ while decaying; saturation; irradiation followed by removal.}
{$\Act=\Prod(1-2^{-t/\Thalf})$ and $N=\Act/\lambda$ during irradiation. After removal at $t_0$:
$\Act=\Act(t_0)\,2^{-(t-t_0)/\Thalf}$. Saturation activity equals $\Prod$.}
{Assuming saturation is reached after one or two half-lives, or forgetting that decay continues after
removal.}
{A sample is irradiated at a constant production rate for two half-lives of the product and then
removed. One half-life after removal, its activity as a fraction of the saturation activity is:}
{$\tfrac34$}
{$\tfrac78$}
{$\tfrac12$}
{$\tfrac38$\archans{D}{At removal $\Act=\Prod(1-\tfrac14)=\tfrac34\Prod$; one half-life later it is
$\tfrac12\times\tfrac34=\tfrac38$ of $\Prod$. Option (A) forgets the decay after removal; (B) carries on
irradiating for a third half-life; (C) assumes saturation was reached before removal.}}

\archq{First-order kinetics in chemistry}
{A reaction ``first order'', or a percentage completed in a given time; asked for another percentage.}
{The same law as decay: the fraction left is $2^{-t/t_{1/2}}$, so a percentage completed fixes the
number of half-lives. $75\%$ complete $=$ two half-lives; $87.5\%$ $=$ three; $99.9\%\approx$ ten.}
{Scaling times linearly with the percentage completed.}
{A first-order reaction is $75\%$ complete in \SI{30}{min}. The time for it to be $87.5\%$ complete is:}
{\SI{35}{min}}
{\SI{60}{min}}
{\SI{45}{min}}
{\SI{15}{min}\archans{C}{$75\%$ complete means $\tfrac14$ left, two half-lives, so $t_{1/2}=\SI{15}{min}$;
$87.5\%$ complete means $\tfrac18$ left, three half-lives, \SI{45}{min}. Option (A) scales linearly;
(B) doubles the time; (D) is the half-life itself.}}

\section{Practice Problems}\label{sec:practice-decaylaw}

The problems below follow the order of the chapter and are meant to be worked without looking back.
Where a constant is needed, use $N_A=\SI{6.022e23}{\per\mole}$ and $\SI{1}{\year}=\SI{3.156e7}{s}$.
Answers, with brief working, follow the list.

\begin{enumerate}[itemsep=3pt]
\item What fraction of a radioactive sample remains after three mean lives?
\item What percentage of a sample has decayed after $1.5$ half-lives?
\item One-sixteenth of a sample remains after \SI{2.0}{h}. Find the half-life.
\item What is the probability that a given nucleus survives its first half-life but decays during the
      second?
\item Free neutrons have a mean life of \SI{878}{s}. A beam of neutrons moves at \SI{2200}{m/s}. What
      fraction of the neutrons decays in a flight of \SI{10}{km}?
\item A nucleus with mean life $\tau=\SI{10}{\year}$ has already existed for \SI{10}{\year}. What is
      its expected remaining lifetime? What is the probability that it survives a further
      \SI{10}{\year}?
\item What mass of cobalt-60 ($\Thalf=\SI{5.27}{\year}$, molar mass \SI{59.93}{g/mol}) has an activity of
      \SI{1.00}{\curie}?
\item An adult body contains about \SI{140}{g} of potassium (molar mass \SI{39.10}{g/mol}), of which
      $0.0117\%$ of the atoms are potassium-40 ($\Thalf=\SI{1.248e9}{\year}$). Find the potassium-40
      activity of the body.
\item Find the activity of \SI{1.00}{g} of strontium-90 ($\Thalf=\SI{28.8}{\year}$, molar mass
      \SI{90.0}{g/mol}), a fission product, in becquerels and in curies.
\item A \SI{5000}{\becquerel} source sits in front of a detector that intercepts $2.0\%$ of the emitted
      radiation and registers $50\%$ of what reaches it. How many counts are recorded in one minute
      (background negligible)?
\item A patient receives iodine-131 ($\Thalf=\SI{8.02}{\day}$) with an initial activity of
      \SI{1.0}{\mega\becquerel}. How many decays occur in the patient's body in total, if none of the
      iodine is excreted?
\item Show that for $t\ll\Thalf$ the number of decays in a time $t$ is approximately $\Act_0t$, and
      estimate the fractional error of this approximation when $t=0.01\,\Thalf$.
\item Raw count rates (including a constant background) are $820$, $420$, $220$ and \SI{120}{\per\second}
      at $t=0$, $1$, $2$ and \SI{3}{h}. With no source present the counter reads \SI{20}{\per\second}.
      Find the half-life. What would one wrongly conclude by ignoring the background?
\item A plot of $\ln\Act$ against $t$ (in days) is a straight line through $(0,\,10.0)$ and
      $(20,\,8.27)$. Find the half-life and the initial activity.
\item Isotope 1 ($\Thalf=\SI{1.0}{h}$) and isotope 2 ($\Thalf=\SI{4.0}{h}$) start with equal numbers of
      nuclei. When are their activities equal?
\item A sample contains two isotopes with half-lives \SI{1.0}{h} and \SI{10}{h}, each initially of
      activity \SI{1000}{\becquerel}. Find the total activity after \SI{5.0}{h}. Which isotope dominates?
\item A semi-log plot of a sample's total activity is curved at early times and straight at late times.
      Explain how the straight part, and peeling (subtracting its extrapolation), reveal both half-lives.
\item Bismuth-212 ($\Thalf=\SI{60.55}{min}$) decays $64.06\%$ by $\beta^-$ emission and $35.94\%$ by
      $\alpha$ emission. Find both partial half-lives.
\item Starting from \num{1.0e6} bismuth-212 nuclei (data as in the previous problem), how many $\alpha$
      particles are emitted (a) in the first \SI{60.55}{min}, (b) in total?
\item Natural uranium ore is billions of years old. Uranium-234 ($\Thalf=\SI{2.455e5}{\year}$) is a
      member of the uranium-238 series ($\Thalf=\SI{4.468e9}{\year}$ for uranium-238). Predict the ratio
      of uranium-234 to uranium-238 atoms in the ore. (The measured value is \num{5.4e-5}.)
\item Old uranium ore contains \num{3.58e-7} radium-226 atoms per uranium-238 atom. Given that radium-226
      has $\Thalf=\SI{1600}{\year}$, deduce the half-life of uranium-238.
\item Radium-226 (\SI{1600}{\year}) decays to radon-222 (\SI{3.82}{\day}). What mass of radon is in
      secular equilibrium with \SI{1.00}{g} of radium in a sealed vessel?
\item A parent with $\Thalf=\SI{4.0}{h}$ feeds a daughter with $\Thalf=\SI{1.0}{h}$. In transient
      equilibrium, what is the ratio of the daughter's activity to the parent's?
\item How many half-lives of irradiation are needed to reach $90\%$ of the saturation activity?
\item A sample is irradiated for two half-lives and then removed. What fraction of the saturation
      activity remains three half-lives after removal?
\item Sodium-24 ($\Thalf=\SI{15}{h}$) is produced in a reactor at a constant rate. Compare the activity
      reached after \SI{15}{h} and after \SI{45}{h} of irradiation, and comment on the cost of the extra
      \SI{30}{h}.
\end{enumerate}

\subsection*{Answers}
{\small
\begin{enumerate}[itemsep=2pt]
\item $e^{-3}=0.050$, that is, $5.0\%$.
\item $1-2^{-1.5}=1-0.354=0.646$, that is, $64.6\%$.
\item $1/16=2^{-4}$: four half-lives in \SI{2.0}{h}, so $\Thalf=\SI{30}{min}$.
\item $e^{-\lambda\Thalf}-e^{-2\lambda\Thalf}=\tfrac12-\tfrac14=\tfrac14$.
\item Flight time $\num{1.0e4}/2200=\SI{4.55}{s}$; fraction $1-e^{-4.55/878}=\num{5.2e-3}$, about
      $0.52\%$.
\item By memorylessness the expected remaining lifetime is still \SI{10}{\year}, and the probability of
      surviving the next \SI{10}{\year} is $e^{-1}=0.37$. Its age is irrelevant.
\item $\lambda=0.693/(5.27\times\num{3.156e7})=\SI{4.17e-9}{\per\second}$; $N=\num{3.7e10}/\lambda=\num{8.9e18}$;
      mass $=(\num{8.9e18}/\num{6.022e23})\times59.93=\SI{8.8e-4}{g}=\SI{0.88}{mg}$.
\item $N=(140/39.10)(\num{6.022e23})(\num{1.17e-4})=\num{2.52e20}$; $\lambda=0.693/(1.248\times10^{9}\times\num{3.156e7})
      =\SI{1.76e-17}{\per\second}$; $\Act=\lambda N\approx\SI{4.4e3}{\becquerel}$: some four thousand
      decays every second inside every adult, from a nuclide made before the Earth formed.
\item $N=\num{6.69e21}$, $\lambda=0.693/(28.8\times\num{3.156e7})=\SI{7.62e-10}{\per\second}$;
      $\Act=\SI{5.1e12}{\becquerel}=\SI{138}{\curie}$, some $30$ times less than a gram of polonium-210:
      the half-life is $76$ times longer, partly offset by the lighter atoms.
\item $5000\times0.020\times0.50\times60=3000$ counts. (This assumes the half-life is much longer than a
      minute, so the activity stays constant during the count.)
\item $N_0=\Act_0/\lambda=\num{1.0e6}\times(8.02\times86\,400)/0.693=\num{1.0e12}$ decays.
\item Exact: $N_0(1-e^{-\lambda t})=(\Act_0/\lambda)(1-e^{-\lambda t})\approx\Act_0t\,(1-\lambda t/2)$, using
      $e^{-x}\approx1-x+x^2/2$. The approximation $\Act_0t$ overestimates by the fraction
      $\lambda t/2=0.693\times0.01/2=0.35\%$.
\item Net rates $800$, $400$, $200$, \SI{100}{\per\second} halve every hour: $\Thalf=\SI{1.0}{h}$.
      Ignoring background, the successive ratios ($1.95$, $1.91$, $1.83$) are not constant: the decay would
      look non-exponential, with an apparently lengthening half-life.
\item Slope $=(8.27-10.0)/20=\SI{-0.0865}{\per\day}=-\lambda$; $\Thalf=0.693/0.0865=\SI{8.0}{\day}$;
      $\Act_0=e^{10.0}=\num{2.2e4}$ (in the plot's activity units).
\item $\lambda_1e^{-\lambda_1t}=\lambda_2e^{-\lambda_2t}$ gives
      $t=\ln(\lambda_1/\lambda_2)/(\lambda_1-\lambda_2)=\ln4/(0.693\times0.75)=\SI{2.67}{h}$. (This is the
      same expression as $t_{\max}$ in a chain: compare Eq.~\eqref{eq:decaylaw-tmax}.)
\item $1000\times2^{-5}+1000\times2^{-0.5}=31+707=\SI{738}{\becquerel}$, dominated by the \SI{10}{h}
      isotope.
\item At late times the short-lived component is negligible, so the straight tail gives the long
      half-life. Extrapolate that line back to $t=0$ and subtract it from the early data; the remainder is
      another straight line on semi-log axes, giving the short half-life.
\item $T_\beta=60.55/0.6406=\SI{94.5}{min}$; $T_\alpha=60.55/0.3594=\SI{168.5}{min}$. Check:
      $1/94.5+1/168.5=1/60.55$.
\item (a) Half the nuclei decay in one half-life, and $35.94\%$ of those emit an $\alpha$ particle:
      \num{1.80e5}. (b) \num{3.59e5}.
\item Secular equilibrium: $N_{234}/N_{238}=T_{234}/T_{238}=\num{2.455e5}/\num{4.468e9}=\num{5.5e-5}$,
      matching the measured value.
\item $T_\text{U}=T_\text{Ra}/\num{3.58e-7}=1600/\num{3.58e-7}=\SI{4.47e9}{\year}$.
\item $m_\text{Rn}=\SI{1.00}{g}\times\dfrac{3.82}{1600\times365.25}\times\dfrac{222}{226}=\SI{6.4e-6}{g}=\SI{6.4}{\micro\gram}$.
\item $\Act_Y/\Act_X=T_X/(T_X-T_Y)=4/3$.
\item $1-2^{-n}=0.9$ gives $n=\log_210=3.32$ half-lives.
\item $\tfrac34\times2^{-3}=\tfrac{3}{32}=0.094$ of saturation.
\item After \SI{15}{h}, $50\%$ of saturation; after \SI{45}{h}, $87.5\%$. Three times the reactor time
      buys only $1.75$ times the activity.
\end{enumerate}
}

\section*{Looking Back on Part A}
\addcontentsline{toc}{section}{Looking Back on Part A}

Part A asked three questions about the nucleus: what it is, how it comes apart, and when. The answers
form a single argument. The nucleus is a tiny, dense drop of protons and neutrons held together by a
short-range force; its mass is less than the mass of its parts, and that missing mass, the binding
energy, is the reservoir of every energy in this volume. Where a rearrangement of nucleons would release
some of that energy, the nucleus is unstable, and the rules of conservation decide which particles come
out and with how much energy. Only the timing escapes prediction for the individual nucleus, and even
the timing is lawful for a sample.

\begin{itemize}
\item \textbf{The nucleus (\chref{ch:nucleus}).} Rutherford scattering, read with the conservation of
      energy and angular momentum, bounds the size of the nucleus; the radius $R=R_0A^{1/3}$ with
      $R_0\approx\SI{1.2}{\femto\metre}$ means every nucleus has the same density. The mass defect gives the
      binding energy $\Eb=\Delta m\,c^2$, and the binding energy per nucleon, peaking near iron and nickel and
      shaped by the terms of the liquid-drop formula, is the map that says which rearrangements release
      energy. The short-range nuclear force and the separation energies complete the picture.
\item \textbf{Radioactive decay (\chref{ch:decay}).} $\alpha$, $\beta^\pm$, electron capture and
      $\gamma$ decay conserve charge, nucleon number, energy, momentum and lepton number; the $Q$-value,
      computed from atomic masses with the electrons counted correctly, is the energy released, shared
      between the products by momentum conservation. The continuous $\beta$ spectrum demanded the
      neutrino. The parabola of isobar masses explains the valley of stability, and $\alpha$ and $\beta$
      steps carry heavy nuclei down the four natural decay series.
\item \textbf{The decay law (this chapter).} A fixed probability per unit time gives
      $N=N_0e^{-\lambda t}$, with $\Thalf=\ln2/\lambda$, $\tau=1/\lambda$ and activity $\Act=\lambda N$,
      measured through a detector's geometry, efficiency and background. Adding isotopes, channels and
      sources to the same equation describes mixtures, branching, chains with their secular and transient
      equilibria, and production at a constant rate; the same equation is first-order chemical kinetics.
\end{itemize}

\noindent\textbf{The thread into Part B.} Part A let nuclei change only of their own accord. Part B makes
them change on demand. \chref{ch:reactions} fires particles at nuclei, applies the same $Q$-value
bookkeeping to reactions, and reads off the binding-energy curve why splitting the heaviest nuclei
(fission) and fusing the lightest (fusion) both release energy. \chref{ch:nucapps} then turns the decay
law of this chapter into a set of tools: the exponential clock dates rocks and bones, the activity and
half-life of a source decide its medical use and the dose it delivers, and a simulation built on nothing
but a fixed probability per step reproduces the law from pure randomness.
